\subsection{PASSAGE FROM LAWS OF INFINITESIMAL OPERATION TO LAWS OF OPERATION}

Lemma 4. Let \(\mathbf{G}\) be a connected topological group, \(\mathbf{X}\) a Hausdorff topological space and \(f_{1}, f_{2}\) laws of left (resp. right) operation of \(\mathbf{G}\) on \(\mathbf{X}\) such that, for all \(x \in \mathbf{X}\) , the mappings \(s \mapsto f_1(s, x)\) , \(s \mapsto f_2(s, x)\) of G into X are continuous. Suppose that there exists a neighbourhood V of \(\{e\} \times X\) in G \(\times\) X such that \(f_1\) and \(f_2\) coincide on V. Then \(f_1 = f_2\) .

Let \(x \in X\) and A be the set of \(g \in G\) such that \(f_{1}(g, x) = f_{2}(g, x)\) . Then A is closed in G. On the other hand, let \(g \in A\) ; we write \(y = f_{1}(g, x) = f_{2}(g, x)\) . There exists a neighbourhood U of e in G such that \(f_{1}(t, y) = f_{2}(t, y)\) for \(t \in U\) , in other words such that \(f_{1}(t', x) = f_{2}(t', x)\) for \(t' \in Ug\) (resp. gU). Hence A is open in G and therefore A = G.

PROPOSITION 16. Let G be a connected Lie group, X a Hausdorff manifold of class \(\mathbf{C}^r\) and \(f_1, f_2\) laws of left (resp. right) operation of class \(\mathbf{C}^r\) of G on X. If the laws of infinitesimal operation associated with \(f_1\) and \(f_2\) are equal, then \(f_1 = f_2\) .

By § 4, no. 7, Corollary to Proposition 11, there exists a neighbourhood V of \(\{e\} \times \mathbf{X}\) in \(\mathbf{G} \times \mathbf{X}\) such that \(f_{1}\) and \(f_{2}\) coincide on V. Hence \(f_{1} = f_{2}\) (Lemma 4).

Lemma 5. Let G be a simply connected topological group, X a Hausdorff topological space, U an open neighbourhood of e in G and \(\psi\) a continuous mapping of U × X into X such that \(\psi(e, x) = x\) and \(\psi(s, \psi(t, x)) = \psi(st, x)\) for all \(x \in X\) and s, t in U such that \(st \in U\) . Let W be a symmetric connected open neighbourhood of e such that \(W^{3} \subset U\) . There exists one and only one law of continuous left operation \(\psi'\) of G on X such that \(\psi'\) and \(\psi\) coincide on W × X. If G is a Lie group and X a manifold of class \(\mathbf{C}^r\) and \(\psi\) is of class \(\mathbf{C}^r\), then \(\psi'\) is of class \(\mathbf{C}^r\).

The uniqueness on \(\psi\) follows from Lemma 4. Let P be the permutation group of X. For \(u \in W^{3}\) , the mapping \(x \mapsto \psi(u, x)\) is an element \(f(u)\) of P and

\[
f \left(u _ {1} u _ {2} u _ {3}\right) = f \left(u _ {1}\right) f \left(u _ {2}\right) f \left(u _ {3}\right)
\]

for \(u_{1}, u_{2}, u_{3}\) in W. Applying Lemma 1 of no. 1, we obtain a morphism \(f'\) of G into P which extends \(f|\mathrm{W}\) . Let \(\psi'(g,x)=f'(g)(x)\) for all \((g,x)\in\mathbf{G}\times\mathbf{X}\) . Then \(\psi'\) is a law of left operation of G on X which coincides with \(\psi\) on W × X. As \(\psi'(g,\psi'(g',x))=\psi'(gg',x)\) for \((g,g',x)\in\mathbf{G}\times\mathbf{G}\times\mathbf{X}\) , the continuity of \(\psi\) on W × X implies the continuity of \(\psi'\) on gW × X for all \(g\in\mathbf{G}\) . Hence \(\psi'\) is continuous. If \(\psi\) is of class \(\mathbf{C}^r\), it is seen similarly that \(\psi'\) is of class \(\mathbf{C}^r\).

THEOREM 5. Let G be a simply connected Lie group, X a compact manifold of class \(\mathbf{C}^r\) ( \(r \geqslant 2\) ) and \(a \mapsto \mathrm{D}_a\) a law of left (resp. right) infinitesimal operation of class \(\mathbf{C}^{r-1}\) of L(G) on X. There exists one and only one law of left (resp. right) operation of class \(\mathbf{C}^{r-1}\) of G on X such that the associated law of infinitesimal operation is \(a \mapsto \mathrm{D}_a\) .

The uniqueness follows from Proposition 16. By § 4, no. 7, Corollary 1 to Theorem 6, there exist a neighbourhood V of \(\{e\} \times \mathbf{X}\) in \(\mathrm{G} \times \mathrm{X}\) and a law chunk of left (resp. right) operation of class \(\mathrm{C}^{r-1}\) of G on X, defined on V, such that the associated law of infinitesimal operation is \(a \mapsto \mathrm{D}_a\) . As X is compact, V can be assumed to be of the form \(U \times X\) , where U is an open neighbourhood of e in G. It then suffices to apply Lemma 5.
% semantic-fragments: legacy-input-seam
\relax{}
